\documentclass[11pt]{article}

%% ============================================================
%% Model memo for scripts/model/tx_gov_trial.R
%% Built from ~/Work/papers/memo-template (see its README for font setup).
%% Compile: pdflatex tx_gov_model.tex  (twice, for refs)
%% ============================================================

\usepackage{amsmath}
\usepackage{mypaperfonts}
\usepackage{microtype}
\usepackage[margin=1in]{geometry}

\usepackage{graphicx}
\usepackage{booktabs}
\usepackage{array}
\newcolumntype{L}[1]{>{\raggedright\arraybackslash}p{#1}}
\usepackage{xcolor}
\definecolor{memolink}{RGB}{0,60,120}
\usepackage[hyphens]{url}
\usepackage[colorlinks=true, linkcolor=memolink, citecolor=memolink,
            urlcolor=memolink]{hyperref}
\usepackage{enumitem}
\setlist{noitemsep, topsep=2pt}

\usepackage{csquotes}

\usepackage[most]{tcolorbox}
\definecolor{memoalert}{RGB}{155,28,49}
\newtcolorbox{memobox}[2][]{%
  enhanced, breakable,
  colback=memoalert!4!white, colframe=memoalert!85!black,
  colbacktitle=memoalert!85!black, coltitle=white,
  fonttitle=\bfseries\small, fontupper=\small,
  title={#2},
  boxrule=0.6pt, arc=1.5pt, titlerule=0pt,
  left=7pt, right=7pt, top=6pt, bottom=6pt,
  before skip=10pt, after skip=10pt,
  #1}

\newcommand{\memoheader}[2]{%
  \noindent{\bfseries #1 \textperiodcentered\ Young-Kyung Kim \textperiodcentered\ #2}%
  \par\vspace{2pt}\hrule\vspace{1.2\baselineskip}%
}

\usepackage{titlesec}
\titleformat{\section}{\normalsize\bfseries}{\thesection.}{0.5em}{}
\titlespacing{\section}{0pt}{1.1\baselineskip}{0.35\baselineskip}
\titleformat{\subsection}{\normalsize\itshape\bfseries}{\thesubsection.}{0.5em}{}
\titlespacing{\subsection}{0pt}{0.9\baselineskip}{0.25\baselineskip}

%% shorthands
\newcommand{\dem}{\mathrm{dem}}
\newcommand{\rep}{\mathrm{rep}}
\newcommand{\oth}{\mathrm{oth}}
\newcommand{\tp}{\mathrm{tp}}
\newcommand{\logit}{\operatorname{logit}}
\newcommand{\Norm}{\mathcal{N}}

\begin{document}

\memoheader{Precinct $\times$ vote-mode election-night model (TX Governor trial)}{October 18, 2026}

\noindent This documents the model implemented in \texttt{scripts/model/tx\_gov\_trial.R}:
what is estimated, how, and what it assumes. The design goal is a forecast of each
candidate's share of all votes that updates as returns arrive by vote mode, with honest
uncertainty under repeated looks. Section~\ref{sec:setup} fixes notation,
\ref{sec:system}--\ref{sec:post} give the model and its estimation,
\ref{sec:sigma}--\ref{sec:agg} the prediction machinery, \ref{sec:int} the intervals,
and \ref{sec:assume} the assumptions and known gaps.

\section{Setup and the composition}\label{sec:setup}

Index precincts $i = 1,\dots,N$ and vote modes $m \in \mathcal{M} = \{\textrm{mail},
\textrm{early}, \textrm{eday}\}$, with $M = |\mathcal{M}| = 3$. The unit of analysis is
the \emph{cell} $(i,m)$. When a cell has reported we observe three counts
$v^{\dem}_{im}, v^{\rep}_{im}, v^{\oth}_{im}$, where $\oth$ collects every candidate
outside the two major parties, including write-ins.

Rather than model the three counts directly, the model works with a factorisation of the
composition:
\begin{equation}
  \tau_{im} = v^{\dem}_{im} + v^{\rep}_{im} + v^{\oth}_{im},
  \qquad
  \tp_{im} = \frac{v^{\oth}_{im}}{\tau_{im}},
  \qquad
  s_{im} = \frac{v^{\dem}_{im}}{v^{\dem}_{im} + v^{\rep}_{im}},
  \label{eq:fwd}
\end{equation}
so $\tau$ is all votes cast in the cell, $\tp$ is the third-party share \emph{of all}
votes, and $s$ is the Democratic share of the \emph{two-party} vote. The map is invertible:
\begin{equation}
  v^{\dem}_{im} = \tau_{im}(1-\tp_{im})s_{im},
  \qquad
  v^{\rep}_{im} = \tau_{im}(1-\tp_{im})(1-s_{im}),
  \qquad
  v^{\oth}_{im} = \tau_{im}\tp_{im}.
  \label{eq:inv}
\end{equation}
Because \eqref{eq:fwd}--\eqref{eq:inv} is a bijection between $(\tau,\tp,s) \in
\mathbb{R}_+ \times (0,1)^2$ and $(v^{\dem},v^{\rep},v^{\oth})$, nothing is lost relative
to modelling the counts, and the three implied shares of all votes sum to one and lie in
$[0,1]$ \emph{by construction} --- no clamping and no post-hoc renormalisation.

This matters because the obvious alternative --- one model per candidate, then divide by
the row sum --- is not innocuous. With two candidates the predictive residuals are drawn
independently when they should be perfectly anticorrelated, and the division cancels part
of the genuine noise; simulation against the implementation puts the resulting predictive
standard deviation about $29\%$ too small. Factorising avoids this and leaves the
inter-candidate dependence to be estimated rather than induced.

\section{The system of equations}\label{sec:system}

Write $k \in \mathcal{K} = \{\tau, \tp, s\}$ for the three outcomes. The precinct-level
covariate vector $x_i$ is deliberately small --- with hundreds of races one cannot
hand-curate a design per contest --- and is given in Table~\ref{tab:covs}. All of the
continuous demographic terms are standardised, so the design mixes a log-count, two rates
and four standardised scores; the $g$-prior of \S\ref{sec:prior} is invariant to that
rescaling, which is why no further normalisation is needed.

\begin{table}[ht]
\centering
\small
\begin{tabular}{@{}L{3.1cm}L{4.3cm}L{6.4cm}@{}}
\toprule
term & source & in which equation \\
\midrule
\texttt{log\_n} & $\log$ registrants, L2 & turnout, third-party, share \\
\texttt{turnout22} & \texttt{vote22}/\texttt{n}, L2 (2022 midterm rate) & turnout, share \\
\texttt{turnout24} & \texttt{vote24}/\texttt{n}, L2 & turnout, share \\
\texttt{age} & median age, voter file & turnout, third-party, share \\
\texttt{male} & male share of registrants & turnout, third-party, share \\
\texttt{pred.bla} & modelled Black share (BISG) & turnout, share \\
\texttt{pred.his} & modelled Hispanic share (BISG) & turnout, share \\
\texttt{dem2} & two-party Dem baseline, \newline \texttt{pr\_harris}/(\texttt{pr\_harris}$+$\texttt{pr\_trump}) & turnout, third-party, share \\
\texttt{tp\_prior\_raw} & $1-$\texttt{pr\_harris}$-$\texttt{pr\_trump} & third-party only \\
\midrule
county indicators & one column per county & all three (penalised; \S\ref{sec:cre}) \\
\bottomrule
\end{tabular}
\caption{The covariate vector $x_i$. Turnout and share share a design; the third-party
equation has its own, smaller one, which is the only non-nested pair and therefore the only
place a SUR efficiency gain is available (\S\ref{sec:assume}).}
\label{tab:covs}
\end{table}

\noindent Note what is \emph{absent}. There is no vote-mode covariate, because mode is a
dimension of the unit rather than a regressor. There is no spatial basis. And the partisan
baseline enters as a single term, \texttt{dem2}, derived from a vendor model calibrated to
2024 precinct returns --- so the design carries one partisan signal, not the three that the
NYC pipeline used. The modelled quantities are
\begin{equation}
  y^{\tau}_{im} = \log(\tau_{im} + 0.1),
  \qquad
  y^{\tp}_{im} = \logit\!\left(\frac{v^{\oth}_{im} + \tfrac12}{\tau_{im} + 1}\right),
  \qquad
  y^{s}_{im}   = \logit\!\left(\frac{v^{\dem}_{im} + \tfrac12}{v^{\dem}_{im}+v^{\rep}_{im} + 1}\right).
  \label{eq:outcomes}
\end{equation}
The $+\tfrac12$ terms are continuity corrections used \emph{only for fitting}: a cell with
zero third-party votes is common and would otherwise give $y^{\tp} = -\infty$.

\paragraph{Turnout is reparameterised as a bounded total times a composition
(2026-10-19).} Equation~\eqref{eq:outcomes} models each cell's turnout as an unconstrained
log-normal, and the three per-mode turnout equations are fitted independently. Nothing then
ties them to each other or to anything physical, and the consequences were severe:
statewide turnout read $184\%$ of truth at 20:00 and $107.8\%$ with $100\%$ of the vote
counted. The prior mode shares are normalised to sum to one, but only in the prior.

The turnout slot now carries a precinct \emph{total} and a \emph{composition}:
\begin{equation}
  \theta_i = \logit\!\left(\frac{T_i}{\kappa\, n_i}\right),
  \qquad
  \eta_{im} = \log\!\frac{p_{im}}{p_{i,\textrm{ref}}}\;\; (m \neq \textrm{ref}),
  \qquad
  \tau_{im} = T_i\,p_{im},
  \quad p_i = \operatorname{softmax}(\eta_i),
  \label{eq:compose}
\end{equation}
so $\sum_m \tau_{im} = T_i$ exactly and $T_i < \kappa n_i$ by construction, with $n_i$ the
registrants and $\kappa = 1.25$ by default. Both transforms are monotone and the outcomes
stay Gaussian, so conjugacy is untouched; the equation count is unchanged at $KM$, one
total plus $M-1$ log-ratios plus the $2M$ share equations, and $\Sigma$ keeps its
dimensions. Two measurements motivate the split. The precinct total is predicted well
($\operatorname{corr}$ of logs $0.962$, sd of log ratio $0.246$) while the mode split is
not (sd $0.698$ for election day, median ratio $1.87$); and a cap only bites on the total,
since per-mode caps must exceed each mode's observed maximum share of registrants and those
sum to $4.66 n_i$ against an actual total of $0.654 n_i$. The reference mode is early
voting --- not because it predicts the others, which it does not
($R^2 = 0.849$ for $\log(\textrm{eday})$ on $\log(\textrm{early}) + \log n$ against $0.849$
for $\log n$ alone), but because it reports first and is the best measured.

Each of the $K M = 9$ outcomes gets its own equation, and they are tied together only
through their errors:
\begin{equation}
  y^{k}_{im} = x_i' \beta^{k}_{m} + \varepsilon^{k}_{im},
  \qquad
  \boldsymbol\varepsilon_i = \big(\varepsilon^{\tau}_{i,\textrm{mail}},\,
  \varepsilon^{\tp}_{i,\textrm{mail}},\, \varepsilon^{s}_{i,\textrm{mail}},\,
  \varepsilon^{\tau}_{i,\textrm{early}},\, \dots\big)'
  \;\sim\; \Norm(\mathbf{0}, \Sigma),
  \label{eq:sur}
\end{equation}
with $\Sigma \in \mathbb{R}^{9\times 9}$ and $\boldsymbol\varepsilon_i$ independent across
$i$. Note that $m$ is an index, not a multiplicand: $\beta^{k}_{m}$ is a separate
$p$-vector per mode, which is algebraically the same as one pooled regression with mode
dummies interacted with every covariate.

This is Zellner's seemingly unrelated regressions (SUR). The equations share no
parameters. $\Sigma$ carries three substantively distinct kinds of dependence:
\begin{enumerate}
  \item turnout $\leftrightarrow$ share \emph{within} a mode;
  \item turnout $\leftrightarrow$ turnout \emph{across} modes (substitution: heavy early
        voting implies a lighter election day);
  \item share $\leftrightarrow$ share \emph{across} modes --- which would let an observed
        mode inform an unobserved mode of the same precinct.
\end{enumerate}

\noindent Measured on North Carolina 2024 (1{,}118 precincts; MEDSL official returns, the
only state with complete precinct $\times$ mode $\times$ party), items 1 and 3 behave very
differently. \emph{Within} mode, $\rho(\tau,s)$ is real and sign-varying: $+0.204$ for mail
($z=6.9$), $+0.052$ for early (not significant), $-0.140$ for election day ($z=-4.7$).
\emph{Across} modes, only mail\,$\leftrightarrow$\,early is distinguishable from zero
($+0.232$, $z=7.9$); mail\,$\leftrightarrow$\,election-day is $-0.050$ ($z=-1.7$) and
early\,$\leftrightarrow$\,election-day is $+0.039$ ($z=1.3$). Cross-mode \emph{turnout}
correlation is also $\approx 0$ (L2, 5{,}362 Texas precincts, $|\rho|<0.2$ in both 2022 and
2024). $\Sigma$ is therefore imposed \textbf{block-diagonal by mode}: the within-mode blocks
are kept because they are measured, the cross-mode blocks dropped because they are not.

\paragraph{Texas 2024 now measures the same quantities, and disagrees with North Carolina
(2026-10-19).} The 2024 Texas scrape yields real precinct $\times$ mode $\times$ party
results, so the blocks that had to be borrowed can be measured directly. On share
\emph{residuals} --- which is what $\Sigma$ holds, and not raw share correlations, which
are dominated by precinct lean and come out at $0.81$--$0.94$:
\[
\begin{array}{lrr}
\textrm{pair} & \textrm{TX 2024} & \textrm{NC 2024}\\
\textrm{mail}\leftrightarrow\textrm{early} & +0.049\ (t=2.0) & +0.232\\
\textrm{mail}\leftrightarrow\textrm{eday}  & -0.117\ (t=-4.9) & -0.050\\
\textrm{early}\leftrightarrow\textrm{eday} & \mathbf{+0.417}\ (t=22.3) & +0.039
\end{array}
\]
The pair North Carolina says is negligible is the strongest in Texas, and it is the one
that would matter most for the red mirage. Within-mode $\rho(\tau,s)$ also differs: TX
gives mail $+0.008$, early $+0.292$, election day $+0.021$. \textbf{The NC-borrowed
cross-mode block is therefore wrong for Texas in exactly the place it is used.} Until a
second state confirms which pattern generalises, the block-diagonal default stands --- but
it should now be read as ``we do not know'' rather than ``we measured it to be zero''.

\paragraph{The mode partisan differential also flips sign between states.} Measured on
TX 2024 with the specification of \S\ref{sec:impl}: mail $a=+0.769$, early $a=+0.087$,
election day $a=+0.082$, with two-party Democratic shares by mode of $67.3\%$, $53.3\%$ and
$55.9\%$. \textbf{Texas election day runs $2.6$pp more Democratic than early voting} ---
the opposite of North Carolina, and the opposite of the red-mirage pattern the mode
machinery was built for. Borrowing NC's $\delta_m$ for Texas is accordingly worth nothing:
on the 2024 backtest it moves mean $|\textrm{err}|$ from $0.60$pp to $0.59$pp, and even an
oracle $\delta_m$ measured on the target election itself only reaches $0.41$pp.

\section{The prior: a $g$-prior placed on fitted values}\label{sec:prior}

Each equation is given a conjugate prior specified in the units of the outcome rather
than on coefficients. Write $\tilde y^{k}_{\cdot m}$ for the vector of prior fitted values
across all $N$ precincts in mode $m$. These come from an external baseline: expected
turnout split across modes, the 2024 two-party lean shifted by a mode-specific offset, and
a shrunk third-party share,
\begin{equation}
  \tilde\tau_{im} = \widehat{\textrm{est}}_i \cdot \pi_{im},
  \qquad
  \tilde s_{im} = \logit^{-1}\!\big(\logit(\textrm{dem2}_i) + \delta_m\big),
  \qquad
  \tilde\tp_{im} = \tp_i,
\end{equation}
where $\pi_{im}$ are the mode turnout shares and $\delta_m$ the mode partisan offsets.

Despite the name, ``$g$-prior'' is not shorthand for ``Gaussian prior.'' It is Zellner's
specific structure $\beta \sim \Norm(\beta_0, g\sigma^2 (X'X)^{-1})$, in which the prior
covariance is proportional to $(X'X)^{-1}$ --- the same matrix appearing in the OLS
sampling variance --- and $g$ is the scalar multiplier. A ridge prior
$\Norm(0, \lambda^{-1}I)$ is equally Gaussian but is \emph{not} a $g$-prior, and the
difference is exactly what produces items 1 and 3 below. Here $g = 1/w_k$.

The prior is implemented as \emph{data augmentation}. Stack all $N$ precincts as
pseudo-observations with response $\tilde y^{k}$ and weight $w_k = n_k / N$, on top of the
$n_{\textrm{obs}}$ reported rows with weight one:
\begin{equation}
  X_a = \begin{bmatrix} X \\ X_{\textrm{obs}} \end{bmatrix},
  \qquad
  y_a = \begin{bmatrix} \tilde y^{k} \\ y^{k}_{\textrm{obs}} \end{bmatrix},
  \qquad
  W_a = \operatorname{diag}\big(w_k \mathbf{1}_N,\ \mathbf{1}_{n_{\textrm{obs}}}\big),
  \label{eq:aug}
\end{equation}
then solve one weighted least-squares problem. The implied prior on the coefficients is
exactly Zellner's $g$-prior,
\begin{equation}
  \beta^{k}_{m} \;\sim\; \Norm\!\Big(\underbrace{(X'X)^{-1}X'\tilde y^{k}}_{\textrm{OLS projection of the prior}},\;\;
  \sigma_k^2 \big(w_k X'X\big)^{-1}\Big),
  \qquad
  \sigma_k^2 \sim \textrm{Scaled-Inv-}\chi^2(\nu_0, s_0^2),
  \label{eq:gprior}
\end{equation}
with $\nu_0 = 4$ and $s_0 = 0.05$. Three consequences are worth stating plainly.

\begin{enumerate}
  \item \textbf{$n_k$ is an effective sample size in precinct-equivalents.} Because the
        pseudo-rows live in the same design space and share the same $\sigma$, setting
        $n_\tau = 3$ really does mean ``three precincts' worth of information.'' The
        prior's \emph{relative} weight then decays as $n_k/(n_k + n_{\textrm{obs}})$ with
        no intervention, which is why $n_k$ is held fixed through the night rather than
        tuned downward as returns arrive. Shrinking it as well would double-discount the
        prior and overstate confidence in early partial data.
  \item \textbf{The prior mean is a projection, so it can be lossy.} Only the component of
        $\tilde y^{k}$ lying in $\operatorname{col}(X)$ survives \eqref{eq:gprior}. Including
        the baseline itself as a covariate makes the projection exact; omitting it
        discards whatever structure the covariates cannot express.
  \item \textbf{The prior precision is proportional to $X'X$.} It is therefore tightest in
        the directions the design is already most informative about, and vanishes in
        near-null directions. Unlike a ridge penalty $\lambda I$, no choice of $n_k$
        improves the conditioning of the design; the prior fixes
        $n_{\textrm{obs}} < p$ early in the night but not collinearity.
\end{enumerate}

\section{The county random effect}\label{sec:cre}

Equation~\eqref{eq:sur} as written has $\boldsymbol\varepsilon_i$ independent across
precincts. That assumption is false, and it is the single most consequential error the
model can make. For a sum of $N$ precinct-level terms with mean pairwise correlation
$\bar\rho$,
\begin{equation}
  \operatorname{Var}\Big(\sum_{i=1}^{N} X_i\Big) = N\sigma^2\big(1 + (N-1)\bar\rho\big),
  \label{eq:rhobar}
\end{equation}
so the standard deviation is inflated by $\sqrt{1+(N-1)\bar\rho}$ relative to
independence. With $N \approx 4{,}500$ a $\bar\rho$ of one \emph{thousandth} roughly
doubles a statewide interval --- while every other dependence examined in this model moves
it by under one percent.

\subsection{What the data say}

Residuals cluster by county, and not marginally. Measuring the county intraclass
correlation of a top-of-ticket statewide share residual (regressed on another statewide
office on the same ballot, vote-weighted, plus log size), on MEDSL official precinct
returns:

\begin{table}[ht]
\centering
\small
\begin{tabular}{@{}lcccc@{}}
\toprule
 & \multicolumn{2}{c}{county ICC} & \multicolumn{2}{c}{implied sd inflation} \\
\cmidrule(lr){2-3}\cmidrule(lr){4-5}
state & 2022 & 2024 & 2022 & 2024 \\
\midrule
\textbf{TX} & \textbf{0.269} & \textbf{0.196} & $3.07\times$ & $2.72\times$ \\
NC & --- & 0.430 & --- & $3.47\times$ \\
GA & 0.295 & --- & $2.37\times$ & --- \\
MI & 0.381 & 0.149 & $4.61\times$ & $2.90\times$ \\
PA & 0.091 & 0.091 & $3.38\times$ & $3.52\times$ \\
WI & 0.114 & 0.144 & $2.37\times$ & $2.66\times$ \\
AZ & 0.089 & 0.227 & $3.19\times$ & $5.12\times$ \\
OH & 0.193 & 0.131 & $4.31\times$ & $3.70\times$ \\
NV & 0.106 & 0.142 & $2.81\times$ & $3.22\times$ \\
\bottomrule
\end{tabular}
\caption{County ICC of the share residual, nine states and two cycles. The ICC ranges
widely but the inflation is tightly clustered, because inflation depends on ICC $\times$
precincts-per-county and states trade off between the two.}
\label{tab:icc}
\end{table}

\noindent Three things follow. The effect is \emph{universal} --- every state-year shows
it. The inflation is \emph{stable} at $2.4\times$ to $5.1\times$, so any precinct model
treating precincts as independent is overconfident by roughly a factor of three
essentially everywhere. And midterm ICCs generally exceed presidential ones, so for a 2026
forecast the 2022 column is the relevant analogue.

The defaults are the Texas measurements, avoiding cross-state borrowing on the parameter
that matters most: $\text{ICC}_s = 0.269$ (MEDSL TX 2022, governor on attorney general,
7{,}529 precincts in 233 counties) and $\text{ICC}_\tau = 0.200$ (L2 Texas,
$\log(\texttt{vote22})$ on registration, demographics and 2018 turnout, 5{,}433 precincts
in 32 counties).

A coarse stratification will not substitute. Regressing the NC county effects on county
attributes gives $R^2 = 0.005$ for lean, $0.281$ for log size, $0.396$ for both with cubic
terms, and $0.034$ for a six-level stratification on lean. Roughly $60\%$ of the county
effect is county-\emph{idiosyncratic}, which is why neither existing pipeline captures it:
the NYC model has no grouping term at all, and the TX primary model's
$(1\,|\,\texttt{stratum\_demo})$ and $(1\,|\,\texttt{geostratum})$ have five and six levels
over 254 counties.

\subsection{Implementation: a ridge-equivalent effect that preserves conjugacy}

Let $C$ be the $N \times J$ matrix of county indicators and $u \in \mathbb{R}^J$ the county
effects. The model becomes
\begin{equation}
  y^{k}_{im} = x_i'\beta^{k}_{m} + u^{k}_{m,c(i)} + \varepsilon^{k}_{im},
  \qquad u^{k}_{m,c} \sim \Norm\big(0,\ \sigma_k^2/\lambda_k\big),
  \qquad \lambda_k = \frac{1 - \text{ICC}_k}{\text{ICC}_k}.
  \label{eq:cre}
\end{equation}
Fitting this with \texttt{lme4} would forfeit conjugacy and the closed-form draws of
\eqref{eq:draws}. Instead, note that a Gaussian random effect is exactly a ridge penalty
on the indicator columns, and that a ridge penalty is representable as pseudo-rows. The
augmented system stacks three blocks and is solved by one weighted QR:
\begin{equation}
  \underbrace{\begin{bmatrix} X & \mathbf{0} \\ X_{\textrm{obs}} & C_{\textrm{obs}} \\
  \mathbf{0} & I_J \end{bmatrix}}_{\textrm{design}},
  \qquad
  \begin{bmatrix} \tilde y^{k} \\ y^{k}_{\textrm{obs}} \\ \mathbf{0}\end{bmatrix},
  \qquad
  W = \operatorname{diag}\big(\tfrac{n_k}{N}\mathbf 1_N,\ w_{\textrm{obs}},\ \lambda_k \mathbf 1_J\big).
  \label{eq:creaug}
\end{equation}
The first block is the $g$-prior of \S\ref{sec:prior}, with the county columns zeroed so
that it governs $\beta$ only; the second is the reported data, weighted by counted ballots;
the third shrinks $u$ toward zero with weight $\lambda_k$. Conjugacy, the closed-form
draws, and the precinct-equivalents reading of $n_k$ all survive.

The behaviour that produces the widening is at the edge: a county with \emph{no} reported
precincts has $u_c$ shrunk to zero but retains the full random-effect variance
$\sigma_k^2/\lambda_k$ in the posterior draws. Unreported counties are therefore
genuinely uncertain rather than silently assumed average.

\subsection{Validation}

Coverage over 14 replications with a fresh truth each (county component at the measured
ICCs), single stage at $70\%$ reported, nominal $95\%$:

\begin{table}[ht]
\centering
\small
\begin{tabular}{@{}lccc@{}}
\toprule
 & coverage & mean width & mean $|$error$|$ \\
\midrule
county effect off & $8/14 = 57\%$ & $0.577$pp & $0.326$pp \\
county effect on  & $13/14 = 93\%$ & $1.331$pp & $0.232$pp \\
\midrule
off, at the TX ICCs & $11/14 = 79\%$ & $0.526$pp & $0.218$pp \\
on, at the TX ICCs  & $13/14 = 93\%$ & $1.012$pp & $0.235$pp \\
\bottomrule
\end{tabular}
\caption{The uncorrected model's nominal $95\%$ intervals cover $57$--$79\%$ of the time.
With the county effect, $93\%$ --- indistinguishable from nominal given $n=14$ (SE on
coverage $\approx 5.8$pp).}
\label{tab:cov}
\end{table}

\noindent Two caveats. The synthetic truth is generated with the same ICCs the model
assumes, so this establishes that the machinery is calibrated \emph{when the ICC is
correct}, not that $0.269$ is correct for 2026. And $n=14$ is thin: the $93\%$ figures are
consistent with roughly $82$--$100\%$, while the $57\%$ and $79\%$ figures are far enough
from nominal to survive that noise.

\paragraph{One component remains unmeasured.} Equation~\eqref{eq:cre} captures the
\emph{county} level. A \emph{statewide} component --- a uniform swing the covariates miss
--- is not identifiable from a single election's residuals, because the fitted intercept
absorbs it by construction. The inflation factors in Table~\ref{tab:icc} are therefore
lower bounds, and recovering the statewide term needs either multi-election backtesting or
an explicit prior, which is what the systematic term in \S\ref{sec:int} is.

\section{Posterior and posterior draws}\label{sec:post}

With \eqref{eq:aug} the posterior is closed-form. Let $R$ be the upper-triangular factor
from the QR decomposition of $W_a^{1/2} X_a$, so $R'R = X_a' W_a X_a$. Then
\begin{equation}
  \hat\beta = \big(X_a' W_a X_a\big)^{-1} X_a' W_a y_a,
  \qquad
  \hat\sigma^2 = \frac{\big\| W_a^{1/2}\big(y_a - X_a\hat\beta\big)\big\|^2}{n_{\textrm{obs}} - p},
\end{equation}
and the variance posterior is updated conjugately,
\begin{equation}
  \nu_\star = \nu_0 + N + n_{\textrm{obs}},
  \qquad
  s_\star^2 = \frac{\nu_0 s_0^2 + (N + n_{\textrm{obs}})\,\hat\sigma^2}{\nu_\star}.
\end{equation}
\paragraph{Corrected 2026-10-19 --- the previous version of this was a bug.}
Earlier drafts described the residual sum of squares as running over the \emph{augmented}
rows while the degrees of freedom used only the reported rows, and defended that as
``crude but directionally correct self-protection''. It is not defensible: the numerator
and the denominator then describe different row sets, and the error is unbounded in both
directions. With $n_{\textrm{obs}}$ small relative to $N$ and $J$, the $g$-prior and ridge
blocks dominate the numerator and $\sigma$ explodes --- the election-day equation fitted on
one county's 15 cells returned $\sigma = 2.19$ on the log-turnout scale, so
$\exp(\mu+\varepsilon)$ produced statewide turnout draws of order $10^{12}$. In the other
direction, setting $\nu_\star = \nu_0 + N + n_{\textrm{obs}}$ counts the pseudo-rows as
data, which makes $\sigma$ effectively \emph{known} and strips its uncertainty out of every
predictive interval; late-night widths collapsed to $0.02$pp while the point estimate was
off by $0.75$pp. Both symptoms came from the same mis-specification.

Residual variance is now computed on the data rows alone, with the degrees of freedom
consumed measured exactly rather than assumed:
\begin{equation}
  \widehat{\mathrm{edf}} = \operatorname{tr}\!\big(H_{\mathcal{D}\mathcal{D}}\big)
    = \sum_{r \in \mathcal{D}} \|Q_{r\cdot}\|^2,
  \qquad
  \nu = n_{\textrm{obs}} - \widehat{\mathrm{edf}},
  \qquad
  \hat\sigma^2 = \frac{\sum_{r\in\mathcal{D}} \big(W_a^{1/2}(y_a - X_a\hat\beta)\big)_r^2}{\nu},
\end{equation}
where $\mathcal{D}$ indexes the reported rows and $Q$ is the thin $Q$ factor of
$W_a^{1/2}X_a$. The trace matters because the county block is penalised: it consumes well
under $J$ degrees of freedom, and when only a few county levels are identified by data at
all, assuming otherwise is badly wrong. The conjugate update then runs on the data degrees
of freedom, $\nu_\star = \nu_0 + \nu$ and
$s_\star^2 = (\nu_0 s_0^2 + \nu\hat\sigma^2)/\nu_\star$.

One consequence is worth stating plainly, because it cost a day to understand. The old
inflated $\sigma$ was accidentally supplying predictive variance for \emph{unobserved}
cells, which the model has no other source for. Removing the bug removed that too, and
synthetic precinct coverage fell from $0.98$ to $0.65$. The missing component is real and
is treated in \S\ref{sec:int}; the old formula was a badly-scaled proxy for it, not a
substitute.

Draws are then taken directly:
\begin{equation}
  \sigma^{(d)} = s_\star \sqrt{\nu_\star / \chi^2_{\nu_\star}},
  \qquad
  \beta^{(d)} = \hat\beta + R^{-1} z^{(d)},
  \qquad
  z^{(d)} \sim \Norm\big(\mathbf{0},\, \sigma^{(d)2} I_p\big),
  \label{eq:draws}
\end{equation}
which is a multivariate $t$ for $\beta$ marginally. We take $D = 500$ draws.

\begin{memobox}{Why Gaussian rather than Poisson counts}
A Poisson or negative-binomial model of candidate counts is the natural alternative and is
what the Texas primary pipeline uses. The Poisson log-likelihood
$\sum_i \{v_i x_i'\beta - e^{x_i'\beta}\}$ is not quadratic in $\beta$, so no
finite-dimensional conjugate family exists --- Gamma conjugacy applies to a free rate
$\lambda_i$, and imposing $\lambda_i = e^{x_i'\beta}$ with $\beta$ shared across precincts
destroys it. One must then optimise or sample, which returns a point estimate rather than a
posterior, and intervals have to be manufactured some other way. Keeping everything
Gaussian preserves \eqref{eq:draws} in closed form, which is what makes refitting many
races every few minutes feasible.
\end{memobox}

\section{$\Sigma$ and conditioning on reported modes}\label{sec:sigma}

$\Sigma$ is treated as an \emph{input}, estimated once from a past election in which every
mode eventually reported, as the complete-case covariance of the $N \times 9$ residual
matrix (with eigenvalue clipping to restore positive-definiteness).

\begin{memobox}{$\Sigma$ cannot be estimated from tonight's returns}
Modes report in sequence --- mail, then early, then election day --- so the number of
precincts with \emph{all} modes counted is near zero until counting is essentially over.
The share\,$\leftrightarrow$\,share block linking mail to election day, which is precisely
the most useful one, is never identified from the current night's data. It has to come
from history. An earlier implementation estimated $\Sigma$ by complete cases from the live
returns and silently fell back to a diagonal matrix at every reporting stage, making the
whole mechanism inert.
\end{memobox}

Partition precinct $i$'s error vector into the components belonging to reported modes
($o$) and unreported modes ($u$). Prediction for the unreported cells then uses the
conditional normal rather than the marginal:
\begin{equation}
  \boldsymbol\varepsilon_{iu} \mid \boldsymbol\varepsilon_{io}
  \;\sim\;
  \Norm\!\Big(
    \Sigma_{uo}\Sigma_{oo}^{-1} \boldsymbol\varepsilon_{io},\;\;
    \Sigma_{uu} - \Sigma_{uo}\Sigma_{oo}^{-1}\Sigma_{ou}
  \Big),
  \qquad
  \varepsilon^{k(d)}_{im} = y^{k}_{im} - x_i'\beta^{k(d)}_{m}.
  \label{eq:cond}
\end{equation}
Concretely: if a precinct's mail vote lands three points more Democratic than predicted,
its election-day prediction shifts by roughly $\rho\cdot 3 \cdot
(\sigma_{\textrm{eday}}/\sigma_{\textrm{mail}})$ points and its variance shrinks by the
Schur complement. Without \eqref{eq:cond}, observing one mode tells you nothing about
another mode of the same precinct beyond what the covariates already said.

Since the observed/unobserved pattern is shared by many precincts, \eqref{eq:cond} is
computed once per distinct pattern --- at most $2^M = 8$ of them.

\section{Prediction and aggregation}\label{sec:agg}

Within each draw $d$: form the latent outcomes from \eqref{eq:sur} using $\beta^{(d)}$ and
the residual draw from \eqref{eq:cond}, invert the transforms of \eqref{eq:outcomes} and
\eqref{eq:compose}, combine each cell's counted votes with its predicted remainder, and
only then aggregate via \eqref{eq:inv}:

\paragraph{Partial counts enter as a mixture, not a substitution (2026-10-19).} A reported
cell used to have its counted votes substituted as if final. Real snapshots are cumulative
\emph{partial} counts: measured on Texas 2024, when an election-day cell first shows votes
the median cell has $5.4\%$ of its final count in, and mail sits at $66\%$ of final at
19:00. Treating those as final understated statewide turnout by $12.5$pp at 22:30 and
biased the share toward the under-counted mode's opponent. The binary alternative --- drop
the cell until it is ``complete'' --- is worse, because it discards a cell that is $99\%$
counted. Writing $C_{im}$ for votes counted so far in the cell and $C_i = \sum_m C_{im}$,
\begin{equation}
  T^{(d)}_i \;\ge\; C_i,
  \qquad
  r^{(d)}_{im} \;=\; \big(T^{(d)}_i - C_i\big)\,
    \frac{\big[T^{(d)}_i p^{(d)}_{im} - C_{im}\big]_+}{\sum_{m'}\big[T^{(d)}_i p^{(d)}_{im'} - C_{im'}\big]_+},
  \qquad
  v^{\dem(d)}_{im} = C^{\dem}_{im} + r^{(d)}_{im}\big(1-\tp^{(d)}_{im}\big)s^{(d)}_{im},
  \label{eq:mixture}
\end{equation}
and analogously for $\rep$ and $\oth$. A counted-out cell has $r \to 0$ and pins itself; a
barely-started cell is mostly predicted. This needs \textbf{no completeness classifier},
which matters because in production the \texttt{stage} field reaches ``complete'' days
late, uses a different vocabulary per state, is downgraded by our own loader, and has been
observed moving \emph{backwards} ($\mathtt{pct\_rep}$ $0.997 \to 0.269 \to 0.364$). The
floor acts on the precinct total rather than per cell: flooring cells independently lets
$\sum_m \tau_{im}$ exceed $T_i$ again and cost $8$pp of turnout accuracy at $99\%$ counted
when it was first written that way.
\begin{equation}
  S^{(d)}_{\dem} = \frac{\sum_{i,m} \tau^{(d)}_{im}\big(1-\tp^{(d)}_{im}\big)s^{(d)}_{im}}
                        {\sum_{i,m} \tau^{(d)}_{im}},
  \qquad
  S^{(d)}_{\oth} = \frac{\sum_{i,m} \tau^{(d)}_{im}\tp^{(d)}_{im}}{\sum_{i,m} \tau^{(d)}_{im}},
  \label{eq:agg}
\end{equation}
and $S^{(d)}_{\rep}$ analogously. The headline quantities are these shares of all votes;
they satisfy $S_{\dem} + S_{\rep} + S_{\oth} = 1$ identically, since the numerators sum to
$\sum \tau$. The two-party share $S_{\dem}/(S_{\dem} + S_{\rep})$ is reported alongside,
because it is what decides the race.

The ordering is not cosmetic. Because $\exp$, $\logit^{-1}$ and the ratio in \eqref{eq:agg}
are all nonlinear, aggregating within each draw and taking quantiles afterwards is not the
same as --- and is the correct alternative to --- forming per-cell intervals and combining
them.

\section{Intervals}\label{sec:int}

Two intervals are reported for the headline quantity, and they answer different questions.

The \textbf{credible interval} summarises the posterior draws and is a fixed-sample
statement: given the model, there is $1-\alpha$ posterior probability that the quantity
lies in the interval, \emph{for one look at the data}.

The \textbf{confidence sequence} is anytime-valid. It is the set of values whose
posterior-to-prior density ratio has not yet crossed $1/\alpha$,
\begin{equation}
  \mathrm{CS}_\alpha = \left\{ y : \frac{p_{\textrm{post}}(y)}{p_{\textrm{prior}}(y)} \le \frac{1}{\alpha} \right\}.
  \label{eq:cs}
\end{equation}
That ratio is a Bayes factor against the hypothesis that the truth is $y$; by Ville's
inequality it is a non-negative martingale with unit expectation under the truth, so it
exceeds $1/\alpha$ with probability at most $\alpha$ \emph{at any stopping time}. Coverage
therefore holds uniformly across the whole night. This matters because the model refits
every few minutes: with on the order of a hundred looks, a freshly computed $95\%$
credible interval has nowhere near $95\%$ joint coverage, while \eqref{eq:cs} does by
construction.

The price is width. For Gaussian prior and posterior with scales $\sigma_0 > \sigma_1$, the
confidence sequence runs roughly $1.2$--$1.6$ times wider than the credible interval, the
ratio shrinking toward one as $\alpha \to 0$ because both widths grow like
$\sqrt{\log(1/\alpha)}$. When the posterior has not concentrated relative to the prior the
ratio in \eqref{eq:cs} can never reach $1/\alpha$ and the interval is $(-\infty,\infty)$ ---
which is the honest answer early in the night, and something a credible interval will never
tell you.

Note that \eqref{eq:cs} requires the prior predictive as its reference measure, and that
this must include a \emph{systematic} component. A prior built only from idiosyncratic
per-cell noise averages out at $1/\sqrt{NM}$ and comes out tighter than the posterior,
which makes the governing quadratic unbounded and returns an infinite interval at every
stage. Real prior uncertainty is about the statewide swing, which is common to all
precincts and does not average away. That scale is now measured rather than assumed:
across $37$ ordered contest pairs in $9$ states, the statewide two-party share moves
$0.22$--$0.26$ logit units between a prior contest and its target, against the
$\sigma/\sqrt{n_{\textrm{prior}}} \approx 0.098$ the $g$-prior implicitly assumed.

\paragraph{A wide prior on the unreported set (2026-10-19).} The cells still to be counted
within a (county, mode) differ from that county's already-counted ones by more than any
existing term can express. Measured on Texas 2024, the logit gap between the remaining and
the counted share has \emph{mean} $\approx 0$ ($-0.10$ to $+0.13$ across stages) and sd
$0.26$--$0.45$, stable under leave-one-county-out ($0.337$--$0.395$ against a full-sample
$0.373$). Mean zero with large stable dispersion is a variance component, not a predictable
offset --- and the predictor route is separately refuted, since within (county, mode)
pre-election covariates explain only $4.6\%$ of election-day lateness and lateness
$\to$ share is $t=-1.4$. The model therefore adds one shared mean-zero offset per
(county, mode) draw to the share of cells that are not counted out,
$\varepsilon^{s}_{im} \mathrel{+}= u_{c(i),m}$ with
$u \sim \Norm(0, \sigma_u^2)$, $\sigma_u = 0.25$. Being mean-zero it must not move the
point estimate, only the width --- which is the sharpest available check that it is wired
correctly. Nothing already in the model covers it: per-cell residuals average over the
unreported set at $1/\sqrt{k}$, giving $\approx 0.02$ on the logit scale, two orders below
the measured $0.3$; and $u_c$ for a (county, mode) is pinned by that county's
\emph{reported} cells, so it cannot express a reported-versus-unreported difference.

\section{Assumptions and known gaps}\label{sec:assume}

\paragraph{Reporting is ignorable given covariates --- now partly measured, and it
largely holds.} A cell enters the likelihood when it has reported, and early reporters are
not a random sample: on Texas 2024 the early vote counted at 18:45 reads $62.95\%$
two-party Democratic against an early-vote final of $53.29\%$, a raw selection of nearly
ten points \emph{within a single mode}. The question is whether the covariates absorb it.
They largely do. \texttt{dem2} alone explains $R^2 = 0.938$ of early-vote logit share and
the full design $0.963$; after fitting, county-mean residuals do \emph{not} significantly
track when a county reported ($\rho = 0.268$, $t = 1.04$, $p = 0.32$, $n = 16$, weighted
residual sd $0.057$ logit). In the backtest the model removes over half the raw skew: the
counted sample is $+1.76$pp at 19:30 and the estimate lands at $+0.79$pp.

Two cautions. Three of the sixteen counties post early-vote totals \emph{before} poll
close and all three carry negative residuals, which drags that correlation toward zero, so
a weak effect may exist that $n=16$ cannot resolve. And the residual $+0.79$pp is better
explained by the prior's level than by geography: \texttt{dem2} is $52.22\%$ vote-weighted
against an actual two-party $54.33\%$, so the prior pulls down while the skewed counted
sample pulls up. That is a prior-level problem, and loosening the level prior is
\emph{not} the fix --- doing so lets the biased early sample pull harder and degrades the
point estimate (max $|\textrm{err}|$ $1.40 \to 2.23$pp).

\paragraph{Errors are independent across precincts.} This is the most consequential
remaining gap. Summing $N$ precincts with independent errors gives aggregate variance
$O(N)$; with correlated errors it is $O(N^2)$, and that difference is what sets the true
width of a statewide band. Earlier drafts of this memo reported that cross-mode conditioning
improves cell-level share prediction by $21\%$ in RMSE and statewide standard deviation by
$10\%$. \textbf{Both figures were artefacts and are withdrawn.} The first came from a
synthetic generator that planted $\rho \approx 0.88$ across modes, so it measured only that
exploiting an assumption helps; the real value is $\approx 0$ for the
mail\,$\leftrightarrow$\,election-day pair. The second compared two $\Sigma$ matrices that
differed in their \emph{diagonal} as well, so it measured an unreconciled variance estimate
rather than correlation structure. Holding the diagonal fixed and varying only the
off-diagonals moves the statewide standard deviation by $0.09\%$; propagating the measured
NC correlations gives $0.09\%$ statewide, $0.21\%$ for a $500$-precinct district and
$0.86\%$ for a $30$-precinct one. Correlation structure is thus second-order at every
aggregate scale, because the residual term contributes only $O(1/\sqrt{k})$ of total
variance while coefficient uncertainty moves every precinct together and does not average
down. A geostratum random effect or explicit spatial term remains the highest-value
addition.

\paragraph{Partial reporting is treated as all-or-nothing.} A cell counts as observed at
\texttt{STAGE\_DONE} and is otherwise unobserved. In reality a half-counted cell is
asymmetric between the equations: a \emph{weighted} observation for the share equations,
but a \emph{censored} one for turnout, since the count so far is a lower bound on
$\tau_{im}$. Treating it as observed would bias turnout downward across nearly every cell
early in the night.

\paragraph{The mode partisan offsets $\delta_m$ are not yet calibrated, and must be
lean-dependent rather than constant.} Measured on NC 2024, mail runs $+0.660$ logit above
its precinct's overall two-party share on average, with a slope of $-0.297$ on the
precinct's own lean ($t=-20.8$); early and election day are near zero on average but also
lean-dependent ($t=17.5$ and $-11.8$). A constant offset is therefore misspecified, not
approximately right. The \emph{likelihood} already accommodates this, since each mode has
its own $\beta^{s}_{m}$ with the partisan baseline in the design; it was the \emph{prior}
that was wrong, and it is now $\delta_{im} = a_m + b_m \logit(\text{dem2}_i)$. Supplying a
correct lean-dependent offset cuts the error at $10\%$ reported from $1.07$ to $0.17$
percentage points in the trial. The coefficients are left at zero for Texas because NC has
no-excuse mail voting while Texas restricts it to voters 65+, disabled or absent, so the
magnitudes should not be transferred; NC establishes that a nonzero, lean-dependent offset
is needed, not its size. Voter-file sources give
mode turnout composition --- for Texas the 2022 midterm split is mail $4.2\%$, early
$63.8\%$, election day $32.0\%$, against $3.5/76.4/20.2$ in 2024, so a presidential-year
analogue would understate election day by about ten points --- but they record how people
voted, never for whom. The offsets therefore require either precinct $\times$ mode
$\times$ party results, or ecological inference crossing precinct mode composition with
precinct partisan outcomes. Until then the trial reproduces the full red-mirage bias on
purpose, and does so visibly.

\paragraph{The SUR system is not fully exploited.} Coefficients are estimated
equation-by-equation and $\Sigma$ is used only at prediction time. Where designs coincide
this costs nothing: with identical regressors the GLS estimator collapses algebraically to
equation-by-equation least squares. Since the turnout and share equations currently share a
design, the only available gain comes through the third-party equation's distinct design
and is worth about $1\%$. Differentiating the turnout and share designs would raise it to
roughly $10\%$, but that requires genuine exclusion restrictions and an iterative
(Gibbs/feasible-GLS) fit, forfeiting the closed form of \eqref{eq:draws}. Given the cost of
refitting many races continuously, that trade is currently unattractive.

\section{Implementation}\label{sec:impl}

Per refit, per race: nine QR decompositions of the augmented design \eqref{eq:aug}, plus at
most eight small solves for \eqref{eq:cond}, each inverting a $\Sigma_{oo}$ of order at most
nine with one Cholesky of the Schur complement. Nothing iterates and nothing has to
converge.

\begin{table}[ht]
\centering
\small
\begin{tabular}{lll}
\toprule
Equation & Quantity & Code \\
\midrule
\eqref{eq:fwd}, \eqref{eq:outcomes} & observed outcomes, corrected and raw & \texttt{fit\_stage} step 4 \\
\eqref{eq:sur}, \eqref{eq:aug}, \eqref{eq:gprior} & nine conjugate fits & \texttt{fit\_stage} step 4a \\
\eqref{eq:draws} & coefficient draws & \texttt{draw\_beta} \\
$\Sigma$ & source and repair & \texttt{fit\_stage} 4b, \texttt{calibrate\_sigma} \\
\eqref{eq:cond} & conditioning by observation pattern & \texttt{fit\_stage} 4c \\
\eqref{eq:inv}, \eqref{eq:agg} & invert, substitute, aggregate & \texttt{fit\_stage} 4d \\
\eqref{eq:cs} & anytime-valid interval & \texttt{conf\_seq}, \texttt{prior\_draws} \\
\bottomrule
\end{tabular}
\end{table}

\noindent The conjugate regression, the credible interval and the confidence sequence come
from the \texttt{enightmodels} package (\texttt{lm\_bayes}, \texttt{cred\_int},
\texttt{conf\_seq}); everything else is in \texttt{tx\_gov\_trial.R}. The trial runs on
synthetic covariates and synthetic returns when the real inputs are absent, so the
plumbing is testable before data arrives; synthetic mode is labelled loudly at run time.

\section{Validation: the Texas 2024 real reporting cadence}\label{sec:valid}

Earlier drafts closed with a request to backtest on a real reporting order rather than a
synthetic one. That has now been done, and it changed the model.

\paragraph{Why the synthetic orders were not enough.} Two reporting orders had been tried,
both invented. A representative order gave $8/8$ coverage and $|\textrm{err}| \le 0.36$pp
--- too benign. An adversarial order in which small rural counties finish first gave $2/7$
coverage and errors to $-16.5$pp. The real Texas curve is neither: all sixteen counties
report \emph{something} within minutes of poll close, so geographic coverage is nearly
complete from the start, but the first batches are $\approx 97\%$ early vote. Texas posts
$47\%$ of its entire vote in the first batch after poll close, and that batch reads
$8.9$pp more Democratic than the final result. Neither invented order produced that, and
the defects it exposed were invisible to both.

\paragraph{Construction.} The panel is rebuilt from 249 raw county scrapes into $98{,}201$
(precinct, mode, timestamp) cells over 16 counties and 4 modes, reproducing an independent
reference exactly (identical key set, zero vote differences). Snapshots are cumulative
\emph{within} a county and counties are scraped independently, so every statewide figure
uses an as-of join --- each county's latest snapshot at or before $t$. Truth is the final
canvass over the same precinct universe: $2{,}417$ precincts, $52.950\%$ Democratic of all
votes, $54.327\%$ two-party. All predictors predate November 2024: \texttt{dem2} and the
third-party prior from the 2020 Senate race, expected turnout from 2020, demographics and
turnout history from 2018 and 2022.

\paragraph{Result.} Against a naive count-what-is-in benchmark that is $+8.94$pp wrong at
18:45, the model holds $\le 1.4$pp at every stage, with mean $|\textrm{err}|$ of
$0.42$--$0.49$pp across configurations, and the interval collapses as the count completes.
\[
\begin{array}{lrrrr}
\textrm{as-of} & \%\ \textrm{votes} & \textrm{naive} & \textrm{model} & \textrm{width}\\
18{:}45 & 16.2 & +8.94 & -1.40 & 6.97\\
19{:}15 & 53.7 & -0.96 & +0.30 & 3.90\\
21{:}00 & 78.2 & -0.65 & +0.16 & 2.45\\
23{:}00 & 89.2 & -0.82 & -0.17 & 2.50\\
03{:}00 & 99.2 & -0.09 & -0.06 & 0.35\\
\textrm{final} & 100.0 & 0.00 & -0.03 & 0.44
\end{array}
\]

\paragraph{What it found.} Three defects that neither synthetic order could reach, each
documented in the section it belongs to: the variance mis-specification of
\S\ref{sec:post}; partial reporting, \S\ref{sec:agg}, which the all-or-nothing synthetic
generator cannot produce at all; and unbounded per-mode turnout, \S\ref{sec:system}.
Turnout at 20:00 improved from $184\%$ of truth to $88\%$ once the composition was
imposed, and from $107.8\%$ to $103.7\%$ at $100\%$ counted.

\paragraph{What it does not establish.} The fifteen stages are nested snapshots of one
election with one truth, so a coverage count across them is not a coverage \emph{rate} ---
it behaves closer to a single observation, and varying the random seed tests Monte Carlo
stability rather than calibration. The honest claim is narrower than the numbers invite:
the fixes removed clear evidence of \emph{under}-coverage, and the point-estimate curve is
good. Demonstrating calibration needs independent replications --- other states, other
elections --- which the multi-state panel work exists to supply. The synthetic path can no
longer serve as that check either: its generator reveals cells all-or-nothing, so it cannot
exercise either partial-reporting fix, and its coverage figures are uninterpretable until
it emits partial counts.

\section{Next steps}

Items 1, 3 and 5 of the previous list are done or superseded; what remains is reordered by
what the Texas 2024 backtest showed actually binds.

\begin{enumerate}
  \item \textbf{Replicate on a second state-election.} This is now the binding item: every
        number in \S\ref{sec:valid} comes from one election, and calibration cannot be
        established from it. North Carolina is the best candidate --- full-state coverage,
        100 counties, mode $\times$ party available --- and its crosswalk is $98.6\%$
        populated, so it is not blocked.
  \item \textbf{Separate level uncertainty from level shrinkage.} The measured statewide
        swing ($0.22$--$0.26$ logit) is the right scale for interval \emph{width} and the
        wrong one for shrinkage: applying it to the prior degrades the point estimate,
        because a looser prior lets the biased early sample pull the level. The likely
        shape is a mean-zero statewide offset on the predictive draw, as in
        \S\ref{sec:int}, rather than a weaker posterior pull. Implemented but defaulted
        off pending that distinction.
  \item \textbf{An atom at the bound for censored turnout.} Right-censoring is implemented
        (EM, conjugacy preserved) and currently \emph{off}: it fixes early turnout
        ($184\% \to 99\%$) and breaks late ($100\% \to 124\%$), because pure Tobit gives a
        finished cell zero probability of sitting exactly at its bound when $\tau = C$
        exactly is the common case. The fix is a mixture likelihood with a discretised
        probability atom, not a different threshold.
  \item \textbf{Cross-precinct correlation} beyond the county effect --- still the term
        that governs statewide width, and still absent.
  \item \textbf{Make the synthetic generator emit partial cumulative counts} and a
        completeness flag, so it can once again serve as a regression check.
  \item Replace the third-party prior with observed past third-party share once precinct
        identifiers are reconciled across vintages.
\end{enumerate}

\end{document}
